% Lösungen Einheit 2 von 3 – Kegel (zusammenbau v0.1)
\einheitenkopf{Einheit 2 von 3 · Kegel}

\erg{18}{a) Durchmesser; $r = 3$ cm}
\erg{19}{a) Höhe \quad b) $V = \frac{1}{3} \cdot \pi \cdot 3^2 \cdot 7 \approx 66{,}0$ cm³ \quad c) $r = 4$ cm; $V = \frac{1}{3} \cdot \pi \cdot 4^2 \cdot 9 \approx 150{,}8$ cm³ \quad d) $V = \frac{1}{3} \cdot \pi \cdot 2{,}5^2 \cdot 6{,}4 \approx 41{,}9$ cm³ \quad e) $V = \frac{1}{3} \cdot \pi \cdot 6^2 \cdot 11 \approx 414{,}7$ cm³ $\approx 415$ ml \quad f) $s = \sqrt{30^2 + 16^2} = \sqrt{1\,156} = 34$ cm \quad g) $h = \sqrt{29^2 - 20^2} = \sqrt{441} = 21$ m \quad h) $M = \pi \cdot 4 \cdot 11 \approx 138{,}2$ cm² \quad i) $G = \pi \cdot 3^2 \approx 28{,}27$ cm², $M = \pi \cdot 3 \cdot 7 \approx 65{,}97$ cm², $O \approx 94{,}2$ cm² \quad j) Grundkreis als flache Ellipse, $4$ cm breit, mit Mittelpunkt; Höhe $5$ cm senkrecht vom Mittelpunkt, gestrichelt; Spitze mit den beiden Enden der Ellipse verbinden \quad k) Kreis: $r = 3$ cm; Kreisausschnitt: Radius $= 8$ cm (die Mantellinie) \quad l) $h = 3 \cdot 150 : (\pi \cdot 4^2) \approx 9{,}0$ cm \quad m) $M = \pi \cdot 3{,}5 \cdot 6{,}2 \approx 68{,}2$ m²; Kosten $\approx 3\,068$ €; Dachhöhe $= \sqrt{6{,}2^2 - 3{,}5^2} \approx 5{,}12$ m; Gesamthöhe $\approx 23{,}1$ m \quad n) Zylinder: $\pi \cdot 3^2 \cdot 5 \approx 141{,}4$ cm³; Kegel: $\frac{1}{3}$ davon $\approx 47{,}1$ cm³; dreimal, weil das Kegelvolumen ein Drittel ist \quad o) Die Aussage trifft nicht zu. Mit dem Radius $r$ des Zylinders und der Höhe $h$: $V_K = \frac{1}{3} \cdot \pi \cdot (6r)^2 \cdot h = 12 \cdot \pi \cdot r^2 \cdot h$, also das Zwölffache des Zylindervolumens $\pi \cdot r^2 \cdot h$.}

19j)

\kegel{2}{5}{$2$ cm}{$5$ cm}{}

\erg{20}{a) Paul hat den Faktor $\frac{1}{3}$ vergessen und wie beim Zylinder gerechnet. Richtig: $V = \frac{1}{3} \cdot \pi \cdot 4^2 \cdot 15 \approx 251{,}3$ cm³ \quad b) Das Kegelvolumen ist ein Drittel von Grundfläche mal Höhe, das Zylindervolumen Grundfläche mal Höhe; drei Kegel füllen also den Zylinder. \quad c) Kegel mit der Spitze unten; oben der Kreis als Ellipse mit dem Radius $2{,}5$ cm; Höhe $11$ cm gestrichelt von der Mitte zur Spitze \quad d) $V = \frac{1}{3} \cdot \pi \cdot 3^2 \cdot 2 \approx 18{,}8$ m³ – mehr als $15$ m³, eine Fuhre reicht nicht.}

20c)

\kegel{1.25}{2.75}{$2{,}5$ cm}{$11$ cm}{}

